% Claude change 10-01: draft related-literature section for the paper.
% Claude change 10-02: revised against the full texts in ~/Downloads/lit_papers
% and Dropbox/Flooding/Literature (Newhouse, Kousky 2010, Xu and Box-Couillard,
% Philippi and Schiller, Hudson et al., Hovekamp and Wagner, Obolensky, Fabian).
% Structure follows the "Related Literature" subsection of Gruber and Solomon
% (2026) and the contribution paragraphs of Wagner (2022). Strand headings
% match the two literature frames in notes/slides_2026-09-21/slides.tex.
% Citation keys are natbib placeholders; a key list is at the bottom.

\section{Related Literature}

This paper sits at the intersection of five literatures: work on subsidized disaster insurance and adaptation, the health-economics literature on insurance coverage of preventive care, the disaster insurance puzzle, the Samaritan's dilemma, and moral hazard in natural disaster insurance markets.

\paragraph{Subsidized disaster insurance and adaptation.}
The closest papers study the National Flood Insurance Program as an insurance market and ask how its pricing interacts with adaptation. \citet{wagner2022} links house-level characteristics to the universe of NFIP contracts and claims in twenty coastal states and documents that willingness to pay for flood insurance is low: fewer than 60 percent of homeowners in high-risk zones insure even though premiums average two thirds of expected payouts. She also shows that houses built under elevation mandates are 25 percent less likely to insure and 31 percent cheaper to cover at a given price, so the market is adversely selected on observable adaptation rather than on private information. \citet{hovekampwagner2023} use the same claims microdata to estimate the return to the minimum elevation standard for new construction and find that the reduction in flood damage is large relative to the cost of building higher. \citet{grubersolomon2026} ask whether the premium subsidy that Risk Rating 2.0 removed was optimal. They estimate a demand elasticity of $-0.32$, a fiscal spillover of \$173 to \$216 per house in ex post aid for each percentage point of coverage, and a mitigation response of 7 to 11 percent among Florida properties facing larger premium increases, and conclude that the optimal subsidy is about 52 percent, close to the pre-reform level. In a different market, \citet{obolensky2024} shows that U.S. crop insurance subsidies flow increasingly to high-risk regions and slow adaptation there, and that subsidies targeted on regional climate trends restore it.

This paper takes the mitigation margin that \citet{grubersolomon2026} treat as a small response to price and makes it a choice in its own right, measured with its own data: elevation permits, FEMA elevation grants, and declared retrofit costs. The planner in our model chooses a premium subsidy and an investment subsidy jointly, so the question is not whether to subsidize insurance but which instrument buys claims reduction at the lower cost in lost coverage. The empirical fact that organizes the model is that, under the subsidized rating in place through 2021, the premium a pre-FIRM household paid did not fall when the house was elevated. The price channel through which risk-based pricing is supposed to reward adaptation was shut, which is why the small mitigation responses in the literature are what one should expect. Where \citet{hovekampwagner2023} evaluate elevation at construction, we evaluate the retrofit of an existing house, which costs an order of magnitude more and is the margin on which the installed housing stock in the floodplain can adapt.

\paragraph{Preventive care in health insurance.}
The structure of the problem is the one health economists study when they ask how generously an insurer should cover prevention. \citet{ellismanning2007} show that an insurer should cover preventive care more generously than treatment because the insured ignore the effect of their own prevention on the premium pool. \citet{newhouse2006} makes the same point for chronic-disease drugs: when the insured share a common pool, a service that lowers future spending is a case for a subsidy rather than an instance of moral hazard, and the case is stronger the longer the insured remain in the pool. Empirically, \citet{chandragrubermcknight2010} find that higher copayments for physician visits and drugs among California retirees raised hospitalizations, concentrated among the chronically ill, with the savings accruing to the supplemental insurer and the hospital costs to Medicare. \citet{starctown2018} show that Part D plans that also bear medical spending cover drugs more generously than stand-alone drug plans, and that the difference is driven by drugs that reduce medical expenditure, because only the integrated plan internalizes the offset. Behind these results is \citet{ehrlichbecker1972}, who distinguish self-insurance, which lowers the size of a loss, from self-protection, which lowers its probability, and show that market insurance is a substitute for the first but can be a complement to the second when the premium responds to protection.

Two features of our setting change the answer. First, the insurer is the federal government, which also pays the ex post aid, so the externality that \citet{chandragrubermcknight2010} and \citet{starctown2018} document across payers is internal to one budget here, and \citeauthor{newhouse2006}'s tenure condition holds trivially because a floodplain house never leaves the pool. Second, elevation is an irreversible capital cost paid once rather than a flow of care, and it reduces damage per flood rather than the probability of flooding, so in the language of \citet{ehrlichbecker1972} it is self-insurance, and any complementarity with market insurance runs only through the premium. Our evidence that the premium did not respond to elevation is therefore the reason a direct subsidy, rather than a price signal, is the instrument to evaluate.

\paragraph{The disaster insurance puzzle.}
A long literature asks why households fail to insure against low-probability, high-loss events even at favorable premiums. \citet{kunreutherpauly2004} rationalize this with a cost of learning the true probability of a rare event, which acts as a threshold below which the household does not bother to evaluate the contract. For floods, the leading explanation of the take-up pattern is that risk perceptions are driven by salience. \citet{gallagher2014} shows that take-up spikes in the year after a flood and decays back to baseline within a decade, and that communities in the same television market that were not flooded respond at a third of the rate. \citet{kousky2010} finds that after the 1993 Mississippi flood, prices of houses in the 500-year floodplain of St.\ Louis County, where no discount had existed before, fell by about 3 percent, while prices in the 100-year floodplain, which already carried a discount, did not move; she attributes the pre-existing discount to the NFIP's maps and purchase requirement. \citet{xuboxcouillard2024} show that Harvey and Irma raised take-up in counties far from the storms in proportion to their Facebook ties to flooded counties, adding about 81,000 policies outside the flooded areas, with the response concentrated in new policies rather than renewals and no sign of adverse selection among the new buyers. \citet{wagner2022} measures the resulting wedge between willingness to pay and expected benefit and shows it is large enough to reverse the sign of the welfare effect of actuarial pricing. We follow her in treating this wedge as a friction that drives choice but not welfare. In our model the friction matters twice: it is why half of high-risk households are uninsured at subsidized prices, and together with the demand elasticity it determines how many insured households a move to risk-based pricing pushes onto disaster aid.

\paragraph{Ex ante insurance and ex post aid.}
A second strand asks whether the expectation of relief depresses insurance demand. \citet{raschky2013} compare Austria, where partial relief is assured by law, with Germany, where full relief is promised ad hoc, and find that the design of the relief program affects private coverage. \citet{kouskyetal2018} match NFIP policies to FEMA individual-assistance grants and estimate that a dollar of average aid lowers average coverage the following year by about six dollars. In agriculture, \citet{deryuginakirwan2018} find that farmers expecting disaster payments insure less and spend less on inputs, with an elasticity of insurance expenditure to expected aid of about 0.2. \citet{philippischiller2024} study the reverse policy: when Baden-W\"urttemberg replaced ad hoc frost relief for winegrowers with a premium subsidy in 2020, announced as the end of relief, insured acreage rose by 13 percentage points per municipality, most where relief had recently been paid, which they read as the subsidy working as a commitment device against future relief as much as a price cut. \citet{grubersolomon2026} estimate the other direction, from coverage to ex post spending, and find that the fiscal spillover is large enough on its own to justify a subsidy. Our model embeds both directions: aid covers a share of uninsured losses, so a household that drops coverage when its premium rises shifts its losses onto the aid budget, and an elevation subsidy that keeps that household insured while lowering its claims reduces outlays on both programs.

\paragraph{Moral hazard in natural disaster insurance.}
Evidence on moral hazard from flood insurance is clearer on the location margin than on the mitigation margin. \citet{browne2019} find that housing permits and starts in Florida were up to 30 percent higher in communities that joined the NFIP, with the increase concentrated in non-coastal counties. \citet{peraltascott2024} use the roll-out of flood maps to show that NFIP entry raised community population by 4 to 5 percent, with a further 5 percent per standard deviation of historical flood risk, mostly because incumbents stayed. \citet{fabian2024} uses the Risk Rating 2.0 premium changes on the universe of parcels in Florida and North Carolina and finds that the reform cut the share of new development in the floodplain by 10 and 27 percent respectively, with development elasticities to premiums of $-0.11$ and $-0.32$. On the mitigation margin the effects are small. \citet{hudson2017} find no evidence of moral hazard in household surveys from Germany and the United States: insurance status is unrelated to two of three German mitigation measures and positively related to the third, and insured U.S. households are no less likely to have prepared. \citet{kousky2019} reviews the evidence and concludes that the effect of insurance on risk reduction is modest and that no study isolates the premium discount for mitigation. \citet{grubersolomon2026} provide the first price-based estimate for elevation and find a response in Florida that they scale down by roughly four fifths for the national housing stock.

We hold location fixed and study existing houses in the high-risk zone, where the relevant adaptation is a retrofit that costs on the order of \$150,000. Our contribution to this strand is to show why the mitigation response to price is small, namely that the premium paid did not reflect elevation under the subsidized schedule, and to ask the question that follows: if price is a weak instrument for adaptation, whether a direct investment subsidy, financed from the same budget, does better.

% ---------------------------------------------------------------------------
% Citation keys used above (natbib). Add to the .bib as needed.
%   wagner2022              Wagner (2022), AEJ: Economic Policy 14(3)
%   hovekampwagner2023      Hovekamp and Wagner (2023), AEA Papers and Proceedings 113
%   grubersolomon2026       Gruber and Solomon (2026), NBER WP 35408
%   obolensky2024           Obolensky (2024), job market paper (crop insurance)
%   ellismanning2007        Ellis and Manning (2007), JHE 26(6)
%   newhouse2006            Newhouse (2006), JHE 25(5)
%   chandragrubermcknight2010  Chandra, Gruber and McKnight (2010), AER 100(1)
%   starctown2018           Starc and Town (2018/2020), NBER WP 21783; REStud 87(6)
%   ehrlichbecker1972       Ehrlich and Becker (1972), JPE 80(4)
%   kunreutherpauly2004     Kunreuther and Pauly (2004), JRU 28(1)
%   gallagher2014           Gallagher (2014), AEJ: Applied 6(3)
%   kousky2010              Kousky (2010), Land Economics 86(3)
%   xuboxcouillard2024      Xu and Box-Couillard (2024), Economic Inquiry 62(3)
%   raschky2013             Raschky, Schwarze, Schwindt and Zahn (2013), ERE 54(2)
%   kouskyetal2018          Kousky, Michel-Kerjan and Raschky (2018), JEEM 87
%   deryuginakirwan2018     Deryugina and Kirwan (2018), Economic Inquiry 56(2)
%   philippischiller2024    Philippi and Schiller (2024), JRI 91(2)
%   browne2019              Browne, Dehring, Eckles and Lastrapes (2019), JRI 86(4)
%   peraltascott2024        Peralta and Scott (2024), JAERE 11(2)
%   fabian2024              Fabian (2024), job market paper
%   hudson2017              Hudson, Botzen, Czajkowski and Kreibich (2017), Land Economics 93(2)
%   kousky2019              Kousky (2019), Annual Review of Resource Economics 11
